\section{What We Expect to Learn}

The empirical program has two headline margins --- behavioural response to the acceleration (entry and composition) and substitution into cash-settled derivatives --- and each can turn out responsive or null, giving four worlds. Each world teaches something different, and each has a known destination.

First, the \emph{responsive-with-substitution} world: campaign entry falls or tilts towards liquid targets in proportion to predetermined required-trading-days exposure, and Item 6 derivative usage rises in exactly the high-exposure targets where the five-business-day deadline binds. This is the intended headline: the acceleration did not shrink activist stakes so much as move them into instruments the final rule chose not to cover, a regulatory-arbitrage result with direct policy teeth, adjudicating the \$810 million per year projection of \textcite{sec2023modernization} against realised behaviour. This world supports the upper tier of outcomes.

Second, the \emph{responsive-without-substitution} world: entry and composition respond to exposure but derivative usage stays flat. Here the reform genuinely changed who campaigns and where, the chilling projection is validated in direction though perhaps not in magnitude, and the paper's contribution is the first realised measurement of how a disclosure deadline reallocates activist attention across the liquidity distribution.

Third, the \emph{null-with-substitution} world: entry is unmoved (consistent with the SEC's own tables showing most filers had already completed their accumulation within the five-business-day window) but derivatives rise in exposed targets. This says the constraint was never the campaign but the form of ownership, and it isolates the dropped Rule 13d-3(e) as the operative policy choice rather than the headline deadline change.

Fourth, the \emph{null-everywhere} world: activists simply comply on every margin. This is the modal ex ante outcome given the regulator's compliance statistics, and it is still a publishable policy evaluation: a pre-registered, well-powered null on a contested rule, with honest minimum detectable effects, adjudicating a large money projection with the realised data the regulator itself has not revisited. That world lands at the honest tier (JCF/RF), and we regard the chilling/substitution worlds as the upside (RFS and above), not the promise.

In every world the second act stands: the event-time-invariant anatomy of activist target gains, decomposing total returns into the filing return $G_{\mathrm{fil}}$, the post-filing run-up $G_{\mathrm{run}}$, and the bid markup $G_{\mathrm{mkp}}$, shows where takeover value reaches target shareholders and disciplines the capitalisation measurement lemma regardless of how the reform margins resolve.

The milestone mapping is fixed. March 2027 is the milestone talk, not the draft: it requires the first-stage facts, the exposure measure, the model skeleton, and an anatomy pilot. The full draft follows in summer 2027, serving the autumn market. Slippage triggers pre-agreed deletions --- chronologies, the signalling block, the full deal-universe anatomy, the pre-2022 backfill, and finally the EDGAR-only fallback --- rather than schedule denial.
